\documentclass[11pt,a4paper]{article}

\def\courseTitle{Industrial Security (Master)}
\def\sectionNumber{02}
\def\sectionTitle{PLC Ladder / Control-Logic Reverse Engineering}
\def\sheetKind{Solutions}

% =====================================================================
% Shared preamble for exercise sheets and solution sheets.
%
% Define BEFORE \input{}-ing this file:
%   \def\courseTitle{...}
%   \def\sectionNumber{...}
%   \def\sectionTitle{...}
%   \def\sheetKind{Exercise Sheet}   or  Solution Sheet
%
% Optional:
%   \def\courseAuthor{...}
%
% The caller must issue \documentclass{...} (typically article) BEFORE
% loading this preamble.
% =====================================================================

\usepackage[utf8]{inputenc}
\usepackage[T1]{fontenc}
\usepackage[english]{babel}
\usepackage[a4paper, margin=2.2cm]{geometry}
\usepackage{graphicx}
\usepackage{listings}
\usepackage{xcolor}
\usepackage{tikz}
\usepackage{booktabs}
\usepackage{array}
\usepackage{enumitem}
\usepackage{titlesec}
\usepackage{fancyhdr}
\usepackage{hyperref}
\usepackage{amsmath}
\usepackage{amssymb}
\usepackage{tcolorbox}
\tcbuselibrary{skins, breakable}
\usepackage{needspace}

% Discourage solo lines split across pages.
\widowpenalty=10000
\clubpenalty=10000
\raggedbottom

% =====================================================================
% Shared TikZ styles for the Industrial Security lecture series.
% Loaded by every slide deck and exercise sheet that uses TikZ.
% =====================================================================

\usetikzlibrary{
  arrows.meta,
  shapes,
  shapes.geometric,
  shapes.symbols,
  shapes.multipart,
  positioning,
  fit,
  calc,
  backgrounds,
  decorations.pathmorphing,
  decorations.pathreplacing,
  decorations.text,
  patterns,
  matrix,
  shadows
}

% --- colour palette --------------------------------------------------
\definecolor{isOT}{HTML}{1F4E79}     % OT (deep blue)
\definecolor{isIT}{HTML}{6F2DA8}     % IT (purple)
\definecolor{isSafety}{HTML}{C00000} % safety red
\definecolor{isNet}{HTML}{2E7D32}    % network green
\definecolor{isAttack}{HTML}{B23A48} % attack red
\definecolor{isAccent}{HTML}{E08E0B} % amber
\definecolor{isMute}{HTML}{6B6B6B}   % grey
\definecolor{isLight}{HTML}{EDF2F8}  % light background
\definecolor{isPLC}{HTML}{2C5282}
\definecolor{isHMI}{HTML}{2F855A}
\definecolor{isSCADA}{HTML}{6B46C1}

% --- generic node styles --------------------------------------------
\tikzset{
  isnode/.style={
    draw, rounded corners=2pt, minimum height=8mm, minimum width=18mm,
    font=\small, align=center, thick
  },
  isbox/.style={isnode, fill=isLight},
  plc/.style={isnode, fill=isPLC!15, draw=isPLC, thick},
  hmi/.style={isnode, fill=isHMI!15, draw=isHMI, thick},
  scada/.style={isnode, fill=isSCADA!15, draw=isSCADA, thick},
  sensor/.style={isnode, fill=isNet!10, draw=isNet, minimum height=6mm, minimum width=14mm, font=\scriptsize},
  actuator/.style={isnode, fill=isAccent!15, draw=isAccent, minimum height=6mm, minimum width=14mm, font=\scriptsize},
  firewall/.style={isnode, fill=isAttack!15, draw=isAttack, thick},
  server/.style={isnode, fill=isMute!15, draw=isMute!80, thick},
  switch/.style={isnode, fill=isLight, draw=black!60, minimum height=6mm, minimum width=14mm, font=\scriptsize},
  workstation/.style={isnode, fill=white, draw=isOT, thick},
  attacker/.style={isnode, fill=isAttack!20, draw=isAttack, thick, font=\small\bfseries},
  inetnode/.style={draw, ellipse, minimum width=24mm, minimum height=10mm, font=\scriptsize, fill=isLight, thick},
  process/.style={isnode, fill=isOT!15, draw=isOT, thick, minimum width=24mm},
  zone/.style={
    draw, rounded corners=4pt, thick, dashed, inner sep=4mm
  },
  conduit/.style={-{Stealth[length=2.5mm]}, thick, isNet!80!black},
  attack/.style={-{Stealth[length=2.5mm]}, thick, isAttack, dashed},
  flow/.style={-{Stealth[length=2.5mm]}, thick, isOT!70!black},
  netline/.style={thick, black!60},
  axisline/.style={-{Stealth[length=2mm]}, thick, black!70},
  tag/.style={font=\scriptsize\itshape, fill=white, inner sep=1pt},
  level/.style={
    draw, fill=isLight, minimum width=0.85\linewidth, minimum height=8mm,
    font=\small, anchor=west
  },
  brace/.style={decorate, decoration={brace, amplitude=4pt, raise=2pt}},
  legendbox/.style={draw, rounded corners=2pt, fill=isLight, inner sep=2mm, font=\scriptsize, align=left}
}

% --- handy macros ----------------------------------------------------
% \plcicon, \hmiicon etc as simple labelled boxes; useful inline.
\newcommand{\plcicon}[1][PLC]{\tikz[baseline=-0.6ex]{\node[plc, minimum width=10mm, minimum height=5mm, font=\scriptsize, inner sep=1pt]{#1};}}
\newcommand{\hmiicon}[1][HMI]{\tikz[baseline=-0.6ex]{\node[hmi, minimum width=10mm, minimum height=5mm, font=\scriptsize, inner sep=1pt]{#1};}}
\newcommand{\fwicon}[1][FW]{\tikz[baseline=-0.6ex]{\node[firewall, minimum width=10mm, minimum height=5mm, font=\scriptsize, inner sep=1pt]{#1};}}


\providecommand{\courseAuthor}{Industrial Security Lecture Series}
\providecommand{\courseInstitute}{Department of Computer Science}
\providecommand{\companyLogo}{logo}

% --- Corporate branding overrides ------------------------------------
% Picks up logo / author / brand colours from the repo-root branding.tex
% when present; falls back to the defaults above otherwise.
\IfFileExists{../../../branding.tex}{% =====================================================================
% Corporate / institutional branding -- single file to customise the
% look of the entire course (slides, notes, exercises, labs).
%
% HOW TO REBRAND IN UNDER A MINUTE:
%   1. Replace `branding/logo.pdf` with your own logo
%      (PDF preferred; PNG and JPG also work -- name it `logo.<ext>`).
%   2. (Optional) change the two colours below to your brand palette.
%   3. (Optional) override the author/institute lines below.
%   4. Re-run `make` at the repo root. That's it.
%
% Everything in this file has sensible defaults; you can delete a line
% to fall back to the built-in defaults, or comment it out.
%
% This file is loaded automatically by the slides, exercise, and lab
% preambles -- you never need to \input it manually.
% =====================================================================

% --- Logo --------------------------------------------------------------
% Path is resolved from each leaf-build directory; the default points at
% the shared `branding/` folder at the repo root. If you put your logo
% somewhere else, set the full or relative path here (without extension):
\renewcommand{\companyLogo}{../../../branding/logo}

% --- Author / institute (appears on the title page footer) ------------
\renewcommand{\courseAuthor}{}
\renewcommand{\courseInstitute}{}

% --- Brand colours ----------------------------------------------------
% slideAccent      = primary brand colour (title bands, accent rule)
% slideSecondary   = secondary brand colour (section badge, highlight)
% Both use HTML hex codes. Keep enough contrast against white text.
\definecolor{slideAccent}{HTML}{00205B}      % deep navy
\definecolor{slideSecondary}{HTML}{00A0DC}   % light blue accent
}{}

% --- listings -------------------------------------------------------
\definecolor{codebg}{rgb}{0.96,0.96,0.96}
\lstset{
  backgroundcolor=\color{codebg},
  basicstyle=\ttfamily\footnotesize,
  breaklines=true,
  frame=single,
  numbers=left,
  numberstyle=\tiny\color{black!50},
  showstringspaces=false,
  tabsize=2
}

% --- header / footer ------------------------------------------------
\pagestyle{fancy}
\fancyhf{}
\renewcommand{\headrulewidth}{0.4pt}
\renewcommand{\footrulewidth}{0pt}
\lhead{\small \courseTitle}
\rhead{\small Section \sectionNumber{} -- \sheetKind}
\cfoot{\thepage}

% --- task / solution boxes -----------------------------------------
% `before={...}` guards every task/solution box with \needspace so the
% box doesn't start with only one or two lines of breathing room at the
% bottom of a page. If less than ~9 lines remain, LaTeX bumps the box
% to the next page; otherwise it starts here and breaks normally if it
% runs long.
% Boxes are `breakable` so very long solutions don't blow the page,
% but `before={\needspace{14\baselineskip}}` pushes them to a fresh
% page whenever fewer than ~14 lines remain. That covers the vast
% majority of single-question splits in practice.
\newtcolorbox{taskbox}[1]{
  enhanced, breakable,
  colback=isLight, colframe=isOT,
  fonttitle=\bfseries\color{white},
  coltitle=white,
  colbacktitle=isOT,
  title={#1},
  arc=2pt, boxrule=0.6pt,
  left=2mm, right=2mm, top=2mm, bottom=2mm,
  before={\par\vspace{2mm}\needspace{14\baselineskip}\par},
  after={\par\vspace{2mm}}
}

\newtcolorbox{solutionbox}{
  enhanced, breakable,
  colback=isNet!5, colframe=isNet,
  fonttitle=\bfseries\color{white},
  coltitle=white,
  colbacktitle=isNet,
  title={Solution},
  arc=2pt, boxrule=0.6pt,
  left=2mm, right=2mm, top=2mm, bottom=2mm,
  before={\par\vspace{2mm}\needspace{14\baselineskip}\par},
  after={\par\vspace{2mm}}
}

\newtcolorbox{hintbox}{
  enhanced, breakable,
  colback=isAccent!10, colframe=isAccent,
  fonttitle=\bfseries\color{white}, coltitle=white, colbacktitle=isAccent,
  title={Hint},
  arc=2pt, boxrule=0.6pt
}

\newtcolorbox{warnbox}{
  enhanced, breakable,
  colback=isAttack!8, colframe=isAttack,
  fonttitle=\bfseries\color{white}, coltitle=white, colbacktitle=isAttack,
  title={Caution},
  arc=2pt, boxrule=0.6pt
}

% --- section formatting --------------------------------------------
\titleformat{\section}{\Large\bfseries\color{isOT}}{\thesection}{0.5em}{}
\titleformat{\subsection}{\large\bfseries\color{isOT!80!black}}{\thesubsection}{0.5em}{}

% --- title block macro ----------------------------------------------
\newcommand{\sheetheader}{%
  \begin{center}
    {\Large\bfseries \courseTitle}\\[0.3em]
    {\large \sheetKind{} -- Section \sectionNumber}\\[0.2em]
    {\large \sectionTitle}\\[0.2em]
    \rule{\linewidth}{0.4pt}
  \end{center}
  \vspace{0.5em}
}


\begin{document}
\sheetheader

\noindent
\textbf{Note.} Model answers are worked in full so they double as a reference. Where a question is
open (the essay, the rewrite), the rubric hint states what must be present for full marks. Accept
defensible alternatives -- mark on reasoning, on whether each claim is anchored to a recovered byte
or line, and on correct ICS terminology, not on exact wording.

\begin{solutionbox}
\textbf{Exercise 2.1 -- Reconstruct control intent from ST.}\\[2pt]
(1)~I/O map, justified by use:
\par\begin{tabular}{p{2.0cm} p{2.4cm} p{1.6cm} p{6.6cm}}
\toprule
\textbf{Address} & \textbf{Name} & \textbf{Dir.} & \textbf{Justification} \\
\midrule
\texttt{\%IX0.0} & Start    & in  & only in the \texttt{(Start OR Motor)} latch term -- momentary start. \\
\texttt{\%IX0.1} & Stop     & in  & ANDed into the latch; dropping it kills the motor (interlock). \\
\texttt{\%IX0.2} & Guard1   & in  & ANDed interlock, same shape as Stop. \\
\texttt{\%IX0.3} & Guard2   & in  & ANDed interlock, same shape as Stop. \\
\texttt{\%QX0.0} & Run lamp / fast & out & \texttt{M0 AND T1.Q}: energised only after the 3\,s delay. \\
\texttt{\%QX0.1} & Start lamp / slow & out & \texttt{M0 AND NOT T1.Q}: energised during the first 3\,s. \\
\bottomrule
\end{tabular}\\[3pt]
(2)~Intent: a motor (e.g.\ a conveyor) seals in on \texttt{Start}, runs while \texttt{Stop} and
both guards stay closed, and drives a two-phase output -- one coil for the first 3~seconds after
start, the other thereafter. This is the classic ``soft-start indication / two-speed staging''
pattern: \texttt{Q0\_1} marks the spin-up window, \texttt{Q0\_0} marks steady run.\\
(3)~\texttt{T1} is an on-delay timer (TON) driven by \texttt{M0}. While the motor is on but
\texttt{T1.Q} is still false (first 3\,s) only \texttt{Q0\_1} is true; once \texttt{T1.Q} latches
true, only \texttt{Q0\_0} is true. They are mutually exclusive because one gates on \texttt{T1.Q}
and the other on \texttt{NOT T1.Q}; before the motor runs at all, both are false (not exclusive --
both off).\\[2pt]
\emph{Rubric hint:} full marks require recognising the TON, the spin-up window, and that the
mutual exclusivity holds only while \texttt{M0} is true.
\end{solutionbox}

\begin{solutionbox}
\textbf{Exercise 2.2 -- Interlocks and the seal-in latch.}\\[2pt]
(1)~The latch is \texttt{(I0\_0 OR M0)}: \texttt{M0} feeds back into its own input. In scan-cycle
terms, once \texttt{Start} (\texttt{I0\_0}) is true for one scan, \texttt{M0} is set; on the next
scan \texttt{Start} may be false, but \texttt{M0} is still true, so \texttt{(I0\_0 OR M0)} stays
true and the rung re-asserts \texttt{M0}. The output image holds the state across scans -- that is
the seal-in.\\
(2)~\texttt{Stop}, \texttt{Guard1}, \texttt{Guard2} default \texttt{TRUE} and are ANDed in, so the
contact is true when the field device is \emph{healthy}; the logic is wired normally-closed (NC).
A cut wire (or lost power, or a stuck-open contact) reads as \texttt{FALSE}, which ANDs the latch
to \texttt{FALSE} and drops the motor. That is fail-safe: any wiring fault stops the machine rather
than silently disabling the interlock. \texttt{Start} is normally-open (NO): it must go true to
energise, and a cut wire merely prevents starting.\\
(3)~Gate the seal-in on all interlocks being healthy, with the AND \emph{outside} the OR:
\begin{lstlisting}
Permit := Stop AND Guard1 AND Guard2 ;
Motor  := (Start OR Motor) AND Permit ;   (* re-arms only via Start *)
\end{lstlisting}
This is already the original's behaviour: when \texttt{Permit} drops, \texttt{Motor} goes false; a
guard re-closing cannot re-energise it because \texttt{Motor} is already false, so a fresh
\texttt{Start} is required.\\[2pt]
\emph{Rubric hint:} accept any latch that does not auto-restart; the danger to call out is putting
an interlock \emph{inside} the OR, which would let the motor self-restart.
\end{solutionbox}

\begin{solutionbox}
\textbf{Exercise 2.3 -- IEC 61131-3 cross-representation.}\\[2pt]
(1)~Ladder Diagram for \texttt{Motor := (Start OR Motor) AND Stop}:
\begin{lstlisting}
   Start      Stop        Motor
 +--| |--+----|/|--... ---( )--+    (Stop drawn NC: -|/|- if NC contact)
 |       |
 | Motor |
 +--| |--+
\end{lstlisting}
Two parallel contacts (\texttt{Start}, \texttt{Motor}) form the OR; in series with \texttt{Stop}
(an NC contact on the fail-safe wiring) they drive the \texttt{Motor} coil.\\
(2)~Instruction List:
\begin{lstlisting}
LD   Start
OR   Motor
AND  Stop
ST   Motor
\end{lstlisting}
(3)~FBD: an \texttt{OR} block with inputs \texttt{Start} and \texttt{Motor}; its output and
\texttt{Stop} feed an \texttt{AND} block; the \texttt{AND} output wires to \texttt{Motor}.\\
(4)~LD/FBD/SFC are graphical: the source of truth is a layout (rungs, wires, step/transition
graphs) serialised inside the project container, and the compiler discards the geometry. ST/IL are
textual and map almost line-for-line onto the compiled accessor calls. Therefore textual languages
are far easier to lift from a compiled artefact: \texttt{POUS.c} reconstructs ST/IL directly,
whereas recovering a faithful LD/FBD diagram requires re-rendering the boolean expression and
inferring a layout that the original tool may not reproduce identically.\\[2pt]
\emph{Rubric hint:} accept any correct LD/IL/FBD; the load-bearing point is part~(4) -- textual
lifts cleanly, graphical needs a re-render.
\end{solutionbox}

\begin{solutionbox}
\textbf{Exercise 2.4 -- Spot the logic bomb.}\\[2pt]
(1)~Branch~A fires when the hidden counter \texttt{n} exceeds \texttt{1500}. \texttt{n} only
increments on scans where \texttt{Trig} (\texttt{\%IX0.4}) is true, so the trigger condition is:
\texttt{\%IX0.4} has been high for more than 1500 scans (cumulative, since the counter is not
reset).\\
(2)~Stuxnet-style tells: (a)~it keys on an \emph{undocumented input} (\texttt{\%IX0.4}) that the
process description never mentions -- read but unjustified; (b)~it is \emph{time/counter-gated}, so
it never fires under normal data and passes a quick functional test; additionally (c)~branch~A
asserts \texttt{Motor := TRUE} unconditionally, \emph{bypassing} every interlock -- an honest
feature would never override Stop and the guards.\\
(3)~At \texttt{T\#20ms} per scan with \texttt{Trig} held high, \texttt{n} increments once per scan,
so it reaches 1501 after $1501 \times 20\,\text{ms} = 30\,020\,\text{ms} \approx 30\,\text{s}$.
Branch~A first fires about \textbf{30~seconds} after pulsing begins.\\
(4)~Branch~A drops \texttt{Stop}, \texttt{Guard1}, \texttt{Guard2} out of the expression entirely
and forces \texttt{Motor := TRUE}. The motor runs even with an open guard or a pressed stop --
the safety interlocks are defeated, so the conveyor can run with a guard door open or refuse to
stop on the E-stop, a personnel-hazard condition.\\[2pt]
\emph{Rubric hint:} the 30~s figure and ``bypasses interlocks'' are the must-haves; full marks also
note that the bomb survives a functional test because the trigger is cumulative and undocumented.
\end{solutionbox}

\begin{solutionbox}
\textbf{Exercise 2.5 -- Scan-cycle impact.}\\[2pt]
(1)~Added cost $= 0.3 + 0.5 + 64 \times 1.1 = 0.8 + 70.4 = 71.2\,\mu$s. New body
$= 40 + 71.2 = 111.2\,\mu$s. As a fraction of the clean body:
$71.2/40 = 178\%$ added (the body is now $\approx 2.78\times$ its original cost). As a fraction of
the 20\,ms period: $111.2\,\mu\text{s} / 20\,000\,\mu\text{s} = 0.56\%$ of the budget
used (added load alone $= 0.36\%$).\\
(2)~Overrun risk appears when total scan work (body $+$ I/O update $+$ task overhead) approaches
the 20\,ms period. Even ignoring overhead, the body would need to grow to $\approx 20\,000\,\mu$s
-- roughly $500\times$ its clean cost -- before the period is threatened. The attacker's
$71.2\,\mu$s is nowhere near that; it is invisible at the period scale.\\
(3)~``Measure the scan time'' is weak because realistic malicious logic adds microseconds against a
millisecond budget -- a $<1\%$ change lost in normal jitter and easily masked. It only catches gross,
clumsy additions (heavy loops, blocking calls). A stronger control is \emph{content} integrity:
continuous block-CRC / hash monitoring of the compiled blocks against a known-good baseline
(Claroty/Dragos/Nozomi-class), which flags \emph{any} byte change regardless of timing cost.\\[2pt]
\emph{Rubric hint:} accept small arithmetic rounding; the conceptual must-have is ``timing is a
weak side-channel; hash the code instead''.
\end{solutionbox}

\begin{solutionbox}
\textbf{Exercise 2.6 -- MatIEC compilation artefacts.}\\[2pt]
(1)~Strip the accessors: \texttt{\_\_GET\_LOCATED(data\_\_->X,)} reads symbol \texttt{X};
\texttt{\_\_SET\_LOCATED(data\_\_->,Y,,expr)} assigns \texttt{Y := expr}. The body becomes:
\begin{lstlisting}
Motor := (((Start OR Motor) AND Stop) AND Guard1) AND Guard2 ;
Lamp  := Motor ;
\end{lstlisting}
which is the original clean conveyor (left-associated AND, so the parentheses are cosmetic).\\
(2)~Division of labour: \texttt{LOCATED\_VARIABLES.h} maps symbolic names to physical
\texttt{\%I}/\texttt{\%Q} addresses and directions (\texttt{START}~$\to$~\texttt{\%IX0.0}, input) --
it tells you \emph{what each wire is} but nothing about the logic. \texttt{POUS.c} carries the
boolean law (how the symbols combine) but uses symbolic names, so without the header you cannot tie
\texttt{START} to a real terminal. You need both: the header for the I/O map, the C body for the
control law.\\
(3)~Editing only the linked ELF makes the \emph{binary} disagree with both the \texttt{.st} source
and \texttt{POUS.c} (which were regenerated from the untouched source). The lab's Step-4 / Step-6
\emph{integrity} path catches it: \texttt{sha256sum core/openplc} diverges from baseline while the
source/C hashes (and the engineering project) are unchanged -- proving the divergence is on the
controller, not in the project. \texttt{nm}/\texttt{strings} on the ELF can then localise the added
symbol or constant.\\[2pt]
\emph{Rubric hint:} part~(3) must name that source and C stay clean while the binary diverges, and
that hashing the runtime artefact (not the project) is what catches it.
\end{solutionbox}

\begin{solutionbox}
\textbf{Exercise 2.7 -- Forensic: what does this block do?}\\[2pt]
(1)~When the operator's set-point \texttt{Speed\_SP} is in the band \texttt{1064}--\texttt{1410},
the block stashes the operator value in \texttt{Saved\_SP}, overwrites the \emph{actual} actuator
command \texttt{Speed\_Cmd} with the fixed value \texttt{1410}, and feeds the HMI echo
\texttt{Speed\_PV\_Echo} the \emph{saved} (original) value. The operator sees the speed they
entered, while the drive actually runs at the forced 1410. Outside the band the block is a
pass-through, so the tamper is invisible unless the set-point is in range.\\
(2)~A range, not a single value, lets the bomb fire across the entire normal operating window of
the target process while staying dormant during commissioning, idle, or out-of-band test values.
It tells you the attacker knew the legitimate set-point range of this specific process -- target
reconnaissance, not a generic payload (the Stuxnet hallmark of process-specific knowledge).\\
(3)~Mapping to the Stuxnet pattern: the pristine main loop still runs the normal control
(\texttt{ELSE} pass-through and the unchanged \texttt{Motor} logic); the \emph{signature check} is
the \texttt{1064}--\texttt{1410} range test; \emph{record-then-replace} is \texttt{Saved\_SP :=
Speed\_SP} followed by \texttt{Speed\_Cmd := 1410}; the \emph{HMI lies} step is
\texttt{Speed\_PV\_Echo := Saved\_SP}, returning the stale value on read-back.\\
(4)~The highest-value tell is \texttt{Speed\_PV\_Echo := Saved\_SP}: an echo value written from a
stashed copy rather than from the real process variable -- written-but-decoupled-from-reality is the
signature of a man-in-the-middle on the operator's view, which is what makes the attack stealthy.
(The fixed \texttt{1410} and the undocumented \texttt{Saved\_SP} are secondary tells.)\\[2pt]
\emph{Rubric hint:} full marks require ``HMI shows old value while drive runs forced value'' and the
mapping onto record-then-replace; recognising the echo write as the key tell earns the top band.
\end{solutionbox}

\begin{solutionbox}
\textbf{Exercise 2.8 -- Essay: why ladder integrity is hard.}\\[2pt]
\emph{Rubric (5~points).}
\begin{enumerate}\setlength{\itemsep}{0pt}
  \item \textbf{Project-vs-runtime asymmetry (1\,pt).} Must state that the PLC executes compiled
        blocks (MC7, \texttt{.L5K}, bytecode, or the linked ELF) while the engineering tool renders
        the project file; an attacker editing only the runtime leaves the project untouched, so the
        tool reports ``no changes''. Stuxnet is the canonical example (Langner 2013;
        Falliere et al., \emph{W32.Stuxnet Dossier}, Symantec 2011).
  \item \textbf{Two technical obstacles (1\,pt each, cap~2).} Any two of: no code-signing / secure
        boot on older firmware; weak, optional, or absent block CRCs (CRC~$\neq$ cryptographic
        hash); proprietary, often unauthenticated upload protocols (S7comm, TriStation, PCOM) so
        read-back tooling is scarce; tight scan-time budgets that forbid heavy on-PLC
        self-verification; retentive memory and online edits that let logic diverge from any
        baseline; tag aliasing that hides semantics (Bool15 = ``open\_valve'').
  \item \textbf{One organisational obstacle (1\,pt).} Downtime windows that make re-validation
        rare; key-switch left in RUN so downloads need no physical presence; engineering-station
        hygiene (the EWS is the single point of trust and is often Internet-adjacent); thin or
        absent vendor support for integrity features on legacy product.
  \item \textbf{Layered controls + residual risk (1\,pt).} Recommend: key-switch in RUN, signed
        program download where supported, version-controlled project files with review/CI,
        continuous block-CRC/hash monitoring against a known-good baseline, and a hardened
        air-gapped EWS with MFA and application allow-listing. Residual risk to acknowledge: legacy
        firmware that cannot sign or attest; a compromised EWS can still re-baseline a monitor; CRC
        is not collision-resistant; and detection still depends on having captured a trustworthy
        baseline before compromise.
\end{enumerate}
\emph{General:} no invented CVEs or advisory IDs; at least one real cited source (Langner 2013;
Symantec 2011; Beresford, BlackHat 2011; SCADA StrangeLove; CISA CoDe16 advisories 2023;
IEC~61131-3:2013).
\end{solutionbox}

\end{document}
